\section{Introduction}

\subsection{Purpose}

Phenopi is an image-based plant phenotyping system for controlled-environment
experiments. It is designed for repeated RGB imaging of plants and for
extracting quantitative traits such as canopy area and growth over time. The
system combines automated image acquisition, image analysis, quality control,
and growth analysis in a modular workflow. It was developed as a low-cost
alternative to commercial phenotyping platforms and can run on inexpensive
hardware such as a Raspberry Pi camera setup. Users do not need extensive
programming experience, but Phenopi is currently command-line based and assumes
basic familiarity with running terminal commands and working with text-based
configuration files.

\subsection{What Phenopi Does}

Phenopi supports a complete basic workflow for image-based plant phenotyping.
During an experiment, a Raspberry Pi captures images at scheduled time points
and stores them using timestamped filenames. After acquisition, the analysis
pipeline processes the images, segments plant canopies, extracts traits, and
writes the results to output tables. The system also produces diagnostic
outputs that help users check image quality and segmentation accuracy before
interpreting the biological results. Growth analysis tools can then be used to
summarize canopy development over time.

\subsection{What This Manual Covers}

This manual explains how to install Phenopi for automated image acquisition,
prepare an experiment, run the scheduler, process captured images, inspect
quality-control outputs, and perform basic growth analysis. The manual focuses
on practical use. It describes the files users need to edit, the commands they
need to run, the outputs they should expect, and the checks they should perform
when something goes wrong. Detailed implementation notes are kept to a minimum.
Where internal behaviour is described, it is included only when it helps users
understand how to operate or troubleshoot the system.
